\documentclass[10pt,a4paper]{amsart}
\usepackage{amsmath,amssymb,amsthm,graphicx}
\usepackage{verbatim}
\newenvironment{reference}{\comment}{\endcomment}
%\newenvironment{reference}{}{}
\newenvironment{slogan}{\comment}{\endcomment}
\newenvironment{history}{\comment}{\endcomment}

% For commutative diagrams we use Xy-pic
\usepackage[all]{xy}

% We use 2cell for 2-commutative diagrams.
\xyoption{2cell}
% Theorem environments.
%
\theoremstyle{plain}
\newtheorem{theorem}[subsection]{Theorem}
\newtheorem{proposition}[subsection]{Proposition}
\newtheorem{lemma}[subsection]{Lemma}

\theoremstyle{definition}
\newtheorem{definition}[subsection]{Definition}
\newtheorem{example}[subsection]{Example}
\newtheorem{exercise}[subsection]{Exercise}
\newtheorem{situation}[subsection]{Situation}

\theoremstyle{remark}
\newtheorem{remark}[subsection]{Remark}
\newtheorem{remarks}[subsection]{Remarks}

\numberwithin{equation}{subsection}

% Macros
%
\def\lim{\mathop{\mathrm{lim}}\nolimits}
\def\colim{\mathop{\mathrm{colim}}\nolimits}
\def\Spec{\mathop{\mathrm{Spec}}}
\def\Hom{\mathop{\mathrm{Hom}}\nolimits}
\def\Ext{\mathop{\mathrm{Ext}}\nolimits}
\def\SheafHom{\mathop{\mathcal{H}\!\mathit{om}}\nolimits}
\def\SheafExt{\mathop{\mathcal{E}\!\mathit{xt}}\nolimits}
\def\Sch{\mathit{Sch}}
\def\Mor{\mathop{\mathrm{Mor}}\nolimits}
\def\Ob{\mathop{\mathrm{Ob}}\nolimits}
\def\Sh{\mathop{\mathit{Sh}}\nolimits}
\def\NL{\mathop{N\!L}\nolimits}
\def\CH{\mathop{\mathrm{CH}}\nolimits}
\def\proetale{{pro\text{-}\acute{e}tale}}
\def\etale{{\acute{e}tale}}
\def\QCoh{\mathit{QCoh}}
\def\Ker{\mathop{\mathrm{Ker}}}
\def\Im{\mathop{\mathrm{Im}}}
\def\Coker{\mathop{\mathrm{Coker}}}
\def\Coim{\mathop{\mathrm{Coim}}}

% Boxtimes
%
\DeclareMathSymbol{\boxtimes}{\mathbin}{AMSa}{"02}

%
% Macros for moduli stacks/spaces
%
\def\QCohstack{\mathcal{QC}\!\mathit{oh}}
\def\Cohstack{\mathcal{C}\!\mathit{oh}}
\def\Spacesstack{\mathcal{S}\!\mathit{paces}}
\def\Quotfunctor{\mathrm{Quot}}
\def\Hilbfunctor{\mathrm{Hilb}}
\def\Curvesstack{\mathcal{C}\!\mathit{urves}}
\def\Polarizedstack{\mathcal{P}\!\mathit{olarized}}
\def\Complexesstack{\mathcal{C}\!\mathit{omplexes}}
% \Pic is the operator that assigns to X its picard group, usage \Pic(X)
% \Picardstack_{X/B} denotes the Picard stack of X over B
% \Picardfunctor_{X/B} denotes the Picard functor of X over B
\def\Pic{\mathop{\mathrm{Pic}}\nolimits}
\def\Picardstack{\mathcal{P}\!\mathit{ic}}
\def\Picardfunctor{\mathrm{Pic}}
\def\Deformationcategory{\mathcal{D}\!\mathit{ef}}
\usepackage{cleveref}
\crefname{originallemma}{original lemma}{original lemmas}
\Crefname{originallemma}{Original lemma}{Original lemmas}
\makeatletter
\newlabel{chow-lemma-chow-ring-projective-bundle}{{02TX}{source}{}{sourcepointer-1}{}}
\newlabel{chow-lemma-chow-ring-projective-bundle@cref}{{[originallemma][0][]02TX}{[1][0][]source}{}{}{}}
\makeatother
\crefname{originalreference}{original reference}{original references}
\Crefname{originalreference}{Original reference}{Original references}
\makeatletter
\newlabel{lemma-cap-projective-bundle}{{02TW}{source}{}{sourcepointer-2}{}}
\newlabel{lemma-cap-projective-bundle@cref}{{[originalreference][0][]02TW}{[1][0][]source}{}{}{}}
\makeatother
\crefname{originalreference}{original reference}{original references}
\Crefname{originalreference}{Original reference}{Original references}
\makeatletter
\newlabel{lemma-exact-sequence-open}{{02RC}{source}{}{sourcepointer-3}{}}
\newlabel{lemma-exact-sequence-open@cref}{{[originalreference][0][]02RC}{[1][0][]source}{}{}{}}
\makeatother
\crefname{originalreference}{original reference}{original references}
\Crefname{originalreference}{Original reference}{Original references}
\makeatletter
\newlabel{lemma-proper-image}{{0EEG}{source}{}{sourcepointer-4}{}}
\newlabel{lemma-proper-image@cref}{{[originalreference][0][]0EEG}{[1][0][]source}{}{}{}}
\makeatother
\crefname{originalreference}{original reference}{original references}
\Crefname{originalreference}{Original reference}{Original references}
\makeatletter
\newlabel{lemma-pullback-affine-fibres-surjective}{{02TT}{source}{}{sourcepointer-5}{}}
\newlabel{lemma-pullback-affine-fibres-surjective@cref}{{[originalreference][0][]02TT}{[1][0][]source}{}{}{}}
\makeatother
\crefname{originalreference}{original reference}{original references}
\Crefname{originalreference}{Original reference}{Original references}
\makeatletter
\newlabel{lemma-restrict-to-open}{{02RX}{source}{}{sourcepointer-6}{}}
\newlabel{lemma-restrict-to-open@cref}{{[originalreference][0][]02RX}{[1][0][]source}{}{}{}}
\makeatother
\crefname{originalreference}{original reference}{original references}
\Crefname{originalreference}{Original reference}{Original references}
\makeatletter
\newlabel{remark-infinite-sums-rational-equivalences}{{02RZ}{source}{}{sourcepointer-7}{}}
\newlabel{remark-infinite-sums-rational-equivalences@cref}{{[originalreference][0][]02RZ}{[1][0][]source}{}{}{}}
\makeatother
\crefname{originalreference}{original reference}{original references}
\Crefname{originalreference}{Original reference}{Original references}
\makeatletter
\newlabel{situation-setup}{{02QL}{source}{}{sourcepointer-8}{}}
\newlabel{situation-setup@cref}{{[originalreference][0][]02QL}{[1][0][]source}{}{}{}}
\makeatother
\title{Projective space bundle formula: the Chow groups of a projective bundle are free over the Chow groups of the base}
\author{The Stacks Project Authors}
\date{stacks-project:spaces-chow.tex:local-snapshot-20261012; GNU FDL 1.2 or later}
\begin{document}
\maketitle
\noindent\textit{Editorial note.} All delivered mathematical text is the author's own, copied from the pinned Stacks Project chapter file spaces-chow.tex; no statement or proof step was rewritten, reordered or removed.\par\medskip
\noindent\textit{Editorial note.} The counted proof is the author's complete argument: an iterated pushforward computation for injectivity and a decomposition plus induction on the dimension of the base for surjectivity, ending with the reduction to the scheme case\par\medskip
\noindent\textit{Editorial note.} The label lemma-chow-ring-projective-bundle names the counted statement, and the proof environment immediately following it in the source is delivered verbatim as its proof.\par\medskip
\noindent\textit{Editorial note.} The counted statement is carried by the author's own lemma environment in the chapter source.\par\medskip
\noindent\textit{Allowed standard prerequisites (editorial).}\begin{itemize}
\item The projection formula and the computation of the pushforward of powers of the first Chern class of the tautological quotient by the projection, which makes the universal combination the identity on the base
\item Chow groups with rational equivalence for locally finite sums of cycles, and the exact sequences relating the Chow groups of a space, an open subspace and the preimage of its complement
\item The reduction of statements about algebraic spaces to schemes along a quasi-compact open, and the scheme case of the projective bundle formula
\end{itemize}\medskip
\ifdefined\hypertarget\hypertarget{sourcepointer-1}{}\fi
\noindent\textit{Source reference.} Original lemma 02TX: sources/stacks-chow/tags.txt, source lines 3712--3712. Prerequisite pointer invoked at proof line(s) 4511; it names a standard fact of a sibling chapter that the counted argument uses, and every step of the counted proof built on it is retained verbatim in the delivered body.\par\medskip
\ifdefined\hypertarget\hypertarget{sourcepointer-2}{}\fi
\noindent\textit{Source reference.} Original reference 02TW: sources/stacks-spaces-chow/pinned-original-source.tex, source lines 4321--4321. Prerequisite pointer invoked at proof line(s) 4409; it names a standard fact of the same chapter that the counted argument uses, and every step of the counted proof built on it is retained verbatim in the delivered body.\par\medskip
\ifdefined\hypertarget\hypertarget{sourcepointer-3}{}\fi
\noindent\textit{Source reference.} Original reference 02RC: sources/stacks-spaces-chow/pinned-original-source.tex, source lines 1306--1306. Prerequisite pointer invoked at proof line(s) 4490; it names a standard fact of the same chapter that the counted argument uses, and every step of the counted proof built on it is retained verbatim in the delivered body.\par\medskip
\ifdefined\hypertarget\hypertarget{sourcepointer-4}{}\fi
\noindent\textit{Source reference.} Original reference 0EEG: sources/stacks-spaces-chow/pinned-original-source.tex, source lines 820--820. Prerequisite pointer invoked at proof line(s) 4438; it names a standard fact of the same chapter that the counted argument uses, and every step of the counted proof built on it is retained verbatim in the delivered body.\par\medskip
\ifdefined\hypertarget\hypertarget{sourcepointer-5}{}\fi
\noindent\textit{Source reference.} Original reference 02TT: sources/stacks-spaces-chow/pinned-original-source.tex, source lines 3948--3948. Prerequisite pointer invoked at proof line(s) 4429; it names a standard fact of the same chapter that the counted argument uses, and every step of the counted proof built on it is retained verbatim in the delivered body.\par\medskip
\ifdefined\hypertarget\hypertarget{sourcepointer-6}{}\fi
\noindent\textit{Source reference.} Original reference 02RX: sources/stacks-spaces-chow/pinned-original-source.tex, source lines 1951--1951. Prerequisite pointer invoked at proof line(s) 4495; it names a standard fact of the same chapter that the counted argument uses, and every step of the counted proof built on it is retained verbatim in the delivered body.\par\medskip
\ifdefined\hypertarget\hypertarget{sourcepointer-7}{}\fi
\noindent\textit{Source reference.} Original reference 02RZ: sources/stacks-spaces-chow/pinned-original-source.tex, source lines 1991--1991. Prerequisite pointer invoked at proof line(s) 4470; it names a standard fact of the same chapter that the counted argument uses, and every step of the counted proof built on it is retained verbatim in the delivered body.\par\medskip
\ifdefined\hypertarget\hypertarget{sourcepointer-8}{}\fi
\noindent\textit{Source reference.} Original reference 02QL: sources/stacks-spaces-chow/pinned-original-source.tex, source lines 264--264. Prerequisite pointer invoked at proof line(s) 4382; it names a standard fact of the same chapter that the counted argument uses, and every step of the counted proof built on it is retained verbatim in the delivered body.\par\medskip
\begin{lemma}[Projective space bundle formula]
\label{lemma-chow-ring-projective-bundle}
Let $(S, \delta)$ be as in Situation \ref{situation-setup}.
Let $X$ be locally of finite type over $S$.
Let $\mathcal{E}$ be a finite locally free $\mathcal{O}_X$-module
$\mathcal{E}$ of rank $r$. Let $(\pi : P \to X, \mathcal{O}_P(1))$
be the projective bundle associated to $\mathcal{E}$.
The map
$$
\bigoplus\nolimits_{i = 0}^{r - 1}
\CH_{k + i}(X)
\longrightarrow
\CH_{k + r - 1}(P),
$$
$$
(\alpha_0, \ldots, \alpha_{r-1})
\longmapsto
\pi^*\alpha_0 +
c_1(\mathcal{O}_P(1)) \cap \pi^*\alpha_1
+ \ldots +
c_1(\mathcal{O}_P(1))^{r - 1} \cap \pi^*\alpha_{r-1}
$$
is an isomorphism.
\end{lemma}

\begin{proof}
Fix $k \in \mathbf{Z}$. We first show the map is injective.
Suppose that $(\alpha_0, \ldots, \alpha_{r - 1})$ is an element
of the left hand side that maps to zero.
By Lemma \ref{lemma-cap-projective-bundle} we see that
$$
0 = \pi_*(\pi^*\alpha_0 +
c_1(\mathcal{O}_P(1)) \cap \pi^*\alpha_1
+ \ldots +
c_1(\mathcal{O}_P(1))^{r - 1} \cap \pi^*\alpha_{r-1})
= \alpha_{r - 1}
$$
Next, we see that
$$
0 = \pi_*(c_1(\mathcal{O}_P(1)) \cap (\pi^*\alpha_0 +
c_1(\mathcal{O}_P(1)) \cap \pi^*\alpha_1
+ \ldots +
c_1(\mathcal{O}_P(1))^{r - 2} \cap \pi^*\alpha_{r - 2}))
= \alpha_{r - 2}
$$
and so on. Hence the map is injective.

\medskip\noindent
To prove the map is surjective, we will argue exactly
as in the proof of Lemma \ref{lemma-pullback-affine-fibres-surjective}
to reduce to the case of schemes.
We urge the reader to skip the proof.

\medskip\noindent
Let $\beta \in \CH_{k + r - 1}(P)$. Write $\beta = \sum m_j[W_j]$ with
$m_j \not = 0$ and $W_j$ pairwise distinct integral closed subspaces of
$\delta$-dimension $k + r$. Then the family $\{W_j\}$
is locally finite in $P$. Let $Z_j \subset X$ be the ``image''
of $W_j$ as in Lemma \ref{lemma-proper-image}. For any quasi-compact open
$U \subset X$ we see that $\pi^{-1}(U) \cap W_j$
is nonempty only for finitely many $j$. Hence the
collection $Z_j$ of images is a locally finite collection
of integral closed subspaces of $X$.

\medskip\noindent
Consider the fibre product diagrams
$$
\xymatrix{
P_j \ar[r] \ar[d]_{\pi_j} & P \ar[d]^\pi \\
Z_j \ar[r] & X
}
$$
Suppose that $[W_j] \in Z_{k + r - 1}(P_j)$
is rationally equivalent to
$$
\pi_j^*\alpha_{j, 0} +
c_1(\mathcal{O}(1)) \cap \pi_j^*\alpha_{j, 1} +
\ldots +
c_1(\mathcal{O}(1))^{r - 1} \cap \pi_j^*\alpha_{j, r - 1}
$$
for some $(k + i)$-cycle $\alpha_{j, i} \in \CH_{k + i}(Z_j)$. Then
$\alpha_i = \sum m_j \beta_{j, i}$ will be a $(k + i)$-cycle on $X$
and
$$
\pi^*\alpha_0 +
c_1(\mathcal{O}(1)) \cap \pi^*\alpha_1 +
\ldots +
c_1(\mathcal{O}(1))^{r - 1} \cap \pi^*\alpha_{r - 1}
$$
will be rationally equivalent to $\beta$ (see
Remark \ref{remark-infinite-sums-rational-equivalences}).
This reduces us to the case $X$ integral, and
$\alpha = [W]$ for some integral closed subscheme
of $P$ dominating $X$. In particular we may
assume that $d = \dim_\delta(X) < \infty$.

\medskip\noindent
Hence we can use induction on $d = \dim_\delta(X)$.
If $d < k$, then $\CH_{k + r - 1}(X) = 0$ and the lemma holds;
this is the base case of the induction.
Consider a nonempty open $U \subset X$.
Suppose that we can show that
$$
\beta|_{\pi^{-1}(U)} =
\pi^*\alpha_0 +
c_1(\mathcal{O}(1)) \cap \pi^*\alpha_1 +
\ldots +
c_1(\mathcal{O}(1))^{r - 1} \cap \pi^*\alpha_{r - 1}
$$
for some $\alpha_i \in Z_{k + i}(U)$.
By Lemma \ref{lemma-exact-sequence-open} we see that
$\alpha_i = \alpha'_i|_U$ for some $\alpha'_i \in Z_{k + i}(X)$.
By the exact sequences
$\CH_{k + i}(\pi^{-1}(X \setminus U)) \to \CH_{k + i}(P) \to
\CH_{k + i}(\pi^{-1}(U))$
of Lemma \ref{lemma-restrict-to-open}
we see that
$$
\beta -
\left(\pi^*\alpha'_0 +
c_1(\mathcal{O}(1)) \cap \pi^*\alpha'_1 +
\ldots +
c_1(\mathcal{O}(1))^{r - 1} \cap \pi^*\alpha'_{r - 1}\right)
$$
comes from a cycle $\beta' \in \CH_{k + r}(\pi^{-1}(X \setminus U))$.
Since $\dim_\delta(X \setminus U) < d$ we win by
induction on $d$.

\medskip\noindent
In particular, by replacing $X$ by a suitable open we may assume
$X$ is a scheme and we have reduced our problem to
Chow Homology, Lemma \ref{chow-lemma-chow-ring-projective-bundle}.
\end{proof}

\end{document}
